\providecommand{\N}{\mathbb{N}}
\providecommand{\tr}{\operatorname{tr}}
\providecommand{\Lap}{\mathcal{L}}
\providecommand{\Nbr}{\mathcal{N}}
\providecommand{\HASH}{\operatorname{HASH}}
\providecommand{\WL}{\textsc{wl}}

\chapter{Grafos: Redes Neuronales de Grafos}
\label{ch:graphs}

Las rejillas y los grupos, objeto de los \Cref{ch:grids,ch:groups}, comparten un
lujo com\'un: un dominio r\'igido y homog\'eneo en el que todo punto se parece a
cualquier otro y el grupo de simetr\'ias act\'ua de forma transitiva. Una rejilla
de p\'ixeles es la misma en todas partes; una esfera tiene el mismo aspecto tras
cualquier rotaci\'on. La mayor\'ia de los datos relacionales de inter\'es
pr\'actico no disfrutan de tal regularidad. Una mol\'ecula, una red social, un
grafo de citas, un mapa de carreteras o una superficie 3D triangulada son una
colecci\'on irregular de entidades ligadas por relaciones por pares, en la que
distintos nodos tienen distinto n\'umero de vecinos y no existe un orden
can\'onico de esos vecinos. El dominio natural de tales datos es un \emph{grafo}.

Este cap\'itulo desarrolla la tercera de las cinco Ges. Comenzamos con la
simetr\'ia que gobierna los grafos ---la \emph{invarianza por permutaciones}--- y
mostramos c\'omo fuerza un c\'omputo local, basado en vecindarios. Ese c\'omputo,
el \emph{paso de mensajes}, resulta ser el mecanismo \'unico que unifica las
arquitecturas de los cap\'itulos anteriores: una red convolucional es paso de
mensajes en una rejilla, la autoatenci\'on es paso de mensajes en un grafo
completo y una red neuronal de grafos es paso de mensajes en un grafo arbitrario.
Examinamos despu\'es los dos linajes cl\'asicos de la convoluci\'on en grafos
---los m\'etodos \emph{espectrales}, construidos sobre el laplaciano del grafo y
una transformada de Fourier en grafos, y los m\'etodos \emph{espaciales},
construidos directamente sobre los vecindarios--- y mostramos c\'omo la
aproximaci\'on polinomial de los filtros espectrales los une. Por \'ultimo
afrontamos cu\'an \emph{expresivos} son estos modelos, exponiendo los l\'imites
fundamentales que impone la jerarqu\'ia de Weisfeiler--Leman y las patolog\'ias
asociadas del sobre-suavizado y la sobre-compresi\'on.

\section{Simetr\'ia de permutaci\'on}
\label{sec:graphs:permutation}

Un grafo $G = (V, E)$ consta de un conjunto finito de v\'ertices (nodos)
$V = \{1, \dots, n\}$ y un conjunto de aristas $E \subseteq V \times V$.
Asociamos a cada nodo un vector de caracter\'isticas y los reunimos en una matriz
$X \in \R^{n \times d}$ cuya $i$-\'esima fila $x_i \in \R^d$ describe el nodo $i$.
La conectividad se codifica mediante una matriz de adyacencia; en grafos
ponderados los pesos registran informaci\'on m\'etrica o relacional.

\begin{definition}[Grafo ponderado y matriz de adyacencia]
\label{def:graphs:weighted}
Un \emph{grafo ponderado} es una terna $G = (V, E, W)$ con una funci\'on de pesos
$W \colon E \to \R$. Su \emph{matriz de adyacencia} $A \in \R^{n \times n}$ tiene
entradas $A_{ij} = W_{ij}$ si $(i,j) \in E$ y $A_{ij} = 0$ en caso contrario; en
grafos no ponderados $A_{ij} \in \{0,1\}$. El grafo es no dirigido cuando
$A = A^\top$. El \emph{vecindario} de un nodo $v$ es
$\Nbr(v) = \{u \in V : (v,u) \in E\}$, y su \emph{grado} es
$\deg(v) = \sum_{u} A_{vu}$.
\end{definition}

El rasgo decisivo de esta representaci\'on es que arrastra un artefacto
arbitrario: el etiquetado de los nodos. Nada intr\'inseco distingue el ``nodo 1''
del ``nodo 7''; renumerar los \'atomos de una mol\'ecula no cambia la mol\'ecula.
Una permutaci\'on de las etiquetas de los nodos es una simetr\'ia de los datos, y
cualquier modelo que construyamos debe respetarla. Una permutaci\'on act\'ua
simult\'aneamente sobre las filas de $X$ y sobre ambos ejes de $A$.

\begin{definition}[Acci\'on de permutaci\'on]
\label{def:graphs:permutation-action}
Sea $\group_n$ el grupo sim\'etrico de permutaciones de $\{1, \dots, n\}$,
representado por matrices de permutaci\'on $P \in \R^{n \times n}$ (un \'unico $1$
por fila y columna, cero en el resto). Una permutaci\'on $P$ act\'ua sobre las
caracter\'isticas de los nodos y sobre la matriz de adyacencia mediante
\[
  X \mapsto P X, \qquad A \mapsto P A P^\top .
\]
\end{definition}

Queremos que un c\'omputo a nivel de nodo $F(X, A)$ conmute con el reetiquetado, y
que un c\'omputo a nivel de grafo $f(X, A)$ ignore por completo el reetiquetado.
Estas son las dos caras de la simetr\'ia introducidas en el \Cref{ch:priors},
especializadas a $\group_n$.

\begin{definition}[Equivarianza e invarianza por permutaciones]
\label{def:graphs:equiv-inv}
Una funci\'on por nodos $F \colon \R^{n \times d} \times \R^{n \times n} \to
\R^{n \times d'}$ es \emph{equivariante por permutaciones} si
\[
  F(P X, P A P^\top) = P\, F(X, A)
  \qquad \text{para todo } P \in \group_n .
\]
Una funci\'on por grafos $f \colon \R^{n \times d} \times \R^{n \times n} \to
\R^{d'}$ es \emph{invariante por permutaciones} si
$f(P X, P A P^\top) = f(X, A)$ para todo $P \in \group_n$.
\end{definition}

Estas condiciones reproducen la equivarianza y la invarianza exigidas a las capas
convolucionales y equivariantes a grupos, pero ahora el grupo relevante es el
grupo sim\'etrico, discreto y finito, en lugar de las traslaciones o las
rotaciones. El esquema del \Cref{ch:priors} se aplica al pie de la letra: apilar
capas equivariantes, que se componen en una aplicaci\'on equivariante, y a\~nadir
una lectura invariante (un \emph{pooling} invariante por permutaciones, como la
suma o la media sobre los nodos) solo al final, cuando se requiere una \'unica
predicci\'on a nivel de grafo.

La restricci\'on es severa. Las \'unicas aplicaciones lineales sobre
caracter\'isticas de nodos que son equivariantes por permutaciones para
\emph{todo} grafo est\'an generadas por dos operadores: la identidad (cada nodo se
transforma a s\'i mismo) y la agregaci\'on por la matriz de unos (cada nodo suma
las caracter\'isticas de todos los nodos). Para obtener una familia m\'as rica y
consciente de la estructura, debemos dejar que la agregaci\'on respete las
aristas: cada nodo puede combinar su propia caracter\'istica con un agregado de
las caracter\'isticas de sus vecinos, donde el agregado es a su vez invariante al
orden en que se enumeran los vecinos. Esta \'unica observaci\'on ---agregar sobre
un vecindario con un operador invariante por permutaciones y luego actualizar---
es la semilla de toda red neuronal de grafos, y a ella nos dirigimos a
continuaci\'on.

\section{Paso de mensajes}
\label{sec:graphs:message-passing}

Una red neuronal de grafos se construye repitiendo una actualizaci\'on local y
equivariante por permutaciones en la que cada nodo refresca su estado a partir de
los estados de sus vecinos. \citet{gilmer2017neural} cristalizaron esta idea en el
marco del \emph{paso de mensajes}, que desde entonces se ha convertido en la
lengua franca del aprendizaje en grafos.

\begin{definition}[Red neuronal de paso de mensajes]
\label{def:graphs:mpnn}
Una \emph{red neuronal de paso de mensajes} (MPNN, por sus siglas en ingl\'es)
opera sobre un grafo $G = (V, E)$ con caracter\'isticas de nodos
$\{h_v^{(\ell)}\}$ y caracter\'isticas opcionales de aristas $\{e_{vu}\}$. La capa
$\ell$ actualiza cada nodo en dos fases:
\begin{align}
  m_v^{(\ell+1)} &= \bigoplus_{u \in \Nbr(v)}
      M_\ell\!\left(h_v^{(\ell)}, h_u^{(\ell)}, e_{vu}\right)
      && \text{(mensaje y agregaci\'on)}, \\
  h_v^{(\ell+1)} &= U_\ell\!\left(h_v^{(\ell)}, m_v^{(\ell+1)}\right)
      && \text{(actualizaci\'on)},
\end{align}
donde $M_\ell$ es una funci\'on de mensaje aprendible, $U_\ell$ una funci\'on de
actualizaci\'on aprendible y $\bigoplus$ un agregador invariante por
permutaciones (suma, media o m\'aximo).
\end{definition}

El agregador es el componente que sostiene todo el edificio: como $\bigoplus$ no
depende del orden de sus argumentos, la capa es equivariante por permutaciones en
el sentido de la \Cref{def:graphs:equiv-inv}, sean cuales sean las funciones de
mensaje y actualizaci\'on. Tras $\ell$ capas cada nodo ha recogido informaci\'on
de su vecindario a $\ell$ saltos, exactamente como apilar capas convolucionales
agranda el campo receptivo. Una predicci\'on a nivel de grafo se produce mediante
un \emph{pooling} invariante final, $h_G = \bigoplus_{v \in V} h_v^{(L)}$,
seguido de un predictor est\'andar.

La \Cref{fig:graphs:message-passing} representa las tres fases para un solo nodo.

\begin{figure}[h]
\centering
\begin{tikzpicture}[>=Stealth, font=\small,
    centernode/.style={circle, draw=gdlred!70, fill=gdlred!20,
        minimum size=22pt, inner sep=0pt, font=\small\bfseries, line width=1.2pt},
    neighnode/.style={circle, draw=gdlblue!60, fill=gdlblue!20,
        minimum size=20pt, inner sep=0pt, font=\small}]
% Etiquetas de panel
\node[font=\bfseries, text=gdlorange!70!black] at (1.5, 4.3) {1. Mensaje};
\node[font=\bfseries, text=gdlpurple!70!black] at (6.6, 4.3) {2. Agregaci\'on};
\node[font=\bfseries, text=gdlgreen!70!black] at (10.6, 4.3) {3. Actualizaci\'on};
% --- Panel 1: mensajes ---
\begin{scope}
  \node[centernode] (vi) at (1.5, 2.0) {$v$};
  \node[neighnode] (v1) at (0.0, 3.5) {$u_1$};
  \node[neighnode] (v2) at (3.0, 3.5) {$u_2$};
  \node[neighnode] (v3) at (3.1, 1.0) {$u_3$};
  \node[neighnode] (v4) at (0.0, 0.5) {$u_4$};
  \foreach \n in {v1,v2,v3,v4}
    \draw[->, line width=1.4pt, gdlorange!70, shorten >=5pt, shorten <=5pt] (\n) -- (vi);
  \node[font=\footnotesize, text=gdlorange!70!black, anchor=north, align=center] at (1.5, -0.2)
    {$m_{u\to v}=M_\ell(h_v,h_u,e_{vu})$};
\end{scope}
\draw[->, line width=1.6pt, gdlgray!70] (4.2, 2.0) -- (5.0, 2.0);
% --- Panel 2: agregaci\'on ---
\begin{scope}[shift={(5.2,0)}]
  \node[circle, draw=gdlpurple!70, fill=gdlpurple!15, minimum size=38pt,
        inner sep=0pt, line width=1.2pt, font=\large] (agg) at (1.4, 2.0) {$\bigoplus$};
  \node[font=\footnotesize, text=gdlpurple!70!black, anchor=north, align=center] at (1.4, 0.2)
    {$m_v=\bigoplus_{u\in\Nbr(v)} m_{u\to v}$};
\end{scope}
\draw[->, line width=1.6pt, gdlgray!70] (8.2, 2.0) -- (9.0, 2.0);
% --- Panel 3: actualizaci\'on ---
\begin{scope}[shift={(9.2,0)}]
  \node[centernode, minimum size=26pt] (vu) at (1.4, 2.0) {$v$};
  \node[font=\footnotesize, text=gdlgreen!70!black, anchor=north, align=center] at (1.4, 0.6)
    {$h_v'=U_\ell(h_v,m_v)$};
\end{scope}
\end{tikzpicture}
\caption[Paso de mensajes en un grafo]{Las tres fases de una capa de paso de
mensajes para un solo nodo $v$. Cada vecino $u$ emite un mensaje $m_{u\to v}$; los
mensajes se combinan mediante un agregador invariante por permutaciones
$\bigoplus$; el estado del nodo se refresca con la funci\'on de actualizaci\'on
$U_\ell$.}
\label{fig:graphs:message-passing}
\end{figure}

\paragraph{Variantes convolucional, atencional y general.}
Dentro de este marco, la elecci\'on de la funci\'on de mensaje define un espectro
de expresividad y coste~\citep{bronstein2021geometric}. En la variante
\emph{convolucional} el mensaje es una ponderaci\'on fija del vecino, determinada
por la estructura, $M_\ell(h_v, h_u) = c_{vu}\, \psi(h_u)$, con coeficientes
escalares $c_{vu}$ que dependen solo del grafo (por ejemplo, de los grados); esta
es la forma que adopta la red convolucional de grafos de \citet{kipf2017semi}. En
la variante \emph{atencional} el coeficiente se aprende a partir de los extremos,
$c_{vu} = a(h_v, h_u)$, como en la red de atenci\'on de grafos de
\citet{velickovic2017graph}. En la variante plenamente \emph{general} el mensaje
es una funci\'on aprendible arbitraria de ambos extremos, como en la
\Cref{def:graphs:mpnn}. El coste y la flexibilidad crecen a lo largo de este
espectro.

\paragraph{Un mecanismo unificador.}
El poder del paso de mensajes radica en que subsume las arquitecturas de los
cap\'itulos anteriores. Lo precisamos como una sucesi\'on de identificaciones; el
\Cref{ch:grids} trat\'o en detalle el caso convolucional.

\begin{proposition}[La convoluci\'on es paso de mensajes en una rejilla]
\label{prop:graphs:cnn-mpnn}
Una capa convolucional sobre la rejilla $\Z^d$ con n\'ucleo $\{W_\delta\}_{\delta
\in \mathcal{K}}$ es una MPNN cuyo grafo es la rejilla con vecindario
$\Nbr(v) = \{u : v - u \in \mathcal{K}\}$, mensaje
$M(h_v, h_u) = W_{v-u}\, h_u$, agregaci\'on por suma y actualizaci\'on
$U(h_v, m_v) = \sigma(m_v + b)$.
\end{proposition}

\begin{proof}
La convoluci\'on discreta calcula
$h_v^{(\ell+1)} = \sigma\!\big(\sum_{\delta \in \mathcal{K}} W_\delta\,
h_{v-\delta}^{(\ell)} + b\big)$. Sustituyendo $u = v - \delta$ se transforma la
suma sobre desplazamientos del n\'ucleo en una suma sobre vecinos de la rejilla,
$\sigma\!\big(\sum_{u \in \Nbr(v)} W_{v-u}\, h_u^{(\ell)} + b\big)$, que es
exactamente la actualizaci\'on MPNN con el mensaje, la agregaci\'on y la
actualizaci\'on indicados.
\end{proof}

La diferencia crucial entre una convoluci\'on y una capa de grafos general es la
\emph{compartici\'on de pesos} $W_{v-u}$: el peso del mensaje depende solo de la
posici\'on relativa $v-u$, que tiene sentido \'unicamente en un dominio con
estructura de traslaci\'on. En una rejilla esto impone equivarianza traslacional
y reduce el n\'umero de par\'ametros de $O(n^2)$ a $O(\abs{\mathcal{K}})$. En un
grafo arbitrario no existe noci\'on de posici\'on relativa, de modo que el peso
del mensaje debe ser en cambio una funci\'on de las caracter\'isticas de los
nodos (atenci\'on) o de invariantes estructurales como el grado (convoluci\'on).
An\'alogamente, la autoatenci\'on es paso de mensajes en el grafo completo $K_n$
con coeficientes dependientes del contenido
$\alpha_{ij} = \operatorname{softmax}_j\big((Q h_i)^\top (K h_j)/\sqrt{d_k}\big)$ y
actualizaci\'on identidad ---la variante atencional anterior aplicada a una
topolog\'ia todos-con-todos---. As\'i, una capa totalmente conectada, una
convoluci\'on, un \emph{transformer} y una red de grafos difieren solo en la
topolog\'ia del grafo y en c\'omo se determina el peso del mensaje; son instancias
de una sola construcci\'on geom\'etrica.

\section{M\'etodos espectrales frente a espaciales}
\label{sec:graphs:spectral-spatial}

El paso de mensajes es una receta \emph{espacial}: act\'ua directamente sobre los
vecindarios. Existe un punto de vista \emph{espectral} complementario que define
la convoluci\'on en un grafo a trav\'es de un an\'alogo de la transformada de
Fourier. Ambos puntos de vista se encuentran, y comprender su relaci\'on aclara
tanto el espacio de dise\~no como la conexi\'on con el procesamiento cl\'asico de
se\~nales.

\subsection{El laplaciano del grafo}
\label{sec:graphs:laplacian}

El operador central del procesamiento de se\~nales en grafos es el laplaciano, la
contrapartida discreta del operador de Laplace--Beltrami en una variedad
(\Cref{ch:geodesics}).

\begin{definition}[Laplaciano del grafo]
\label{def:graphs:laplacian}
Para un grafo no dirigido y ponderado con matriz de adyacencia $A$ y matriz de
grados $D = \operatorname{diag}(\deg(1), \dots, \deg(n))$, el laplaciano
(combinatorio) es $L = D - A$, y el \emph{laplaciano normalizado} es
\[
  \Lap = D^{-1/2} L\, D^{-1/2} = I - D^{-1/2} A\, D^{-1/2}.
\]
\end{definition}

\begin{figure}[h]
\centering
\begin{tikzpicture}[font=\small,
    gnode/.style={circle, draw=gdlblue!70, fill=gdlblue!20, minimum size=20pt,
        inner sep=0pt, font=\small\bfseries}]
\node[gnode] (1) at (0,1.6) {1};
\node[gnode] (2) at (1.6,2.2) {2};
\node[gnode] (3) at (2.2,0.6) {3};
\node[gnode] (4) at (0.6,0.0) {4};
\draw[gdlgray!60, thick] (1)--(2) (1)--(4) (2)--(3) (3)--(4) (1)--(3);
\node[anchor=west, align=left] at (3.6,1.1)
  {$L = D - A =\begin{bmatrix}
     3 & -1 & -1 & -1\\
     -1 & 2 & -1 & 0\\
     -1 & -1 & 3 & -1\\
     -1 & 0 & -1 & 2
   \end{bmatrix}$};
\end{tikzpicture}
\caption[El laplaciano del grafo]{Un grafo peque\~no y su laplaciano $L = D - A$.
La diagonal registra los grados, $L_{ii} = \deg(i)$, y las entradas fuera de la
diagonal registran la adyacencia, $L_{ij} = -W_{ij}$ para $(i,j) \in E$.}
\label{fig:graphs:laplacian}
\end{figure}

El laplaciano es un operador diferencial discreto: actuando sobre una se\~nal de
nodos $f \colon V \to \R$,
\[
  (L f)(i) = \sum_{j \sim i} W_{ij}\,\big(f(i) - f(j)\big),
\]
la diferencia entre $f(i)$ y un promedio ponderado de sus vecinos ---el an\'alogo
discreto de $-\Delta f = -\operatorname{div}\operatorname{grad} f$---. Sus
propiedades algebraicas y espectrales se resumen a continuaci\'on.

\begin{proposition}[Propiedades del laplaciano]
\label{prop:graphs:laplacian-props}
El laplaciano $L = D - A$ de un grafo no dirigido y ponderado es sim\'etrico y
semidefinido positivo, con
\[
  f^\top L f = \tfrac12 \sum_{(i,j) \in E} W_{ij}\,\big(f(i) - f(j)\big)^2 \ge 0 .
\]
En consecuencia sus valores propios son reales y no negativos,
$0 = \lambda_1 \le \lambda_2 \le \cdots \le \lambda_n$, el vector constante
$\mathbf{1}$ es un vector propio para $\lambda_1 = 0$, y la multiplicidad del
valor propio $0$ es igual al n\'umero de componentes conexas de $G$.
\end{proposition}

\begin{proof}
La simetr\'ia es inmediata, pues $D$ es diagonal y $A = A^\top$. Para la forma
cuadr\'atica, $f^\top L f = \sum_i \deg(i) f(i)^2 - \sum_{i,j} A_{ij} f(i) f(j) =
\tfrac12 \sum_{i,j} A_{ij}(f(i)-f(j))^2$, manifiestamente no negativa; luego
$L \succeq 0$ y todos los $\lambda_i \ge 0$. Como cada fila de $L$ suma cero,
$L \mathbf{1} = 0$, de modo que $\lambda_1 = 0$. Por \'ultimo, $f^\top L f = 0$
obliga a $f(i) = f(j)$ en cada arista, es decir, $f$ es constante en cada
componente conexa; el espacio de tales $f$ tiene dimensi\'on igual al n\'umero de
componentes.
\end{proof}

La forma cuadr\'atica $f^\top L f$ mide la \emph{suavidad} de $f$ sobre el grafo:
es peque\~na cuando $f$ var\'ia poco entre aristas y grande cuando $f$ cambia
abruptamente. Esta es la \emph{energ\'ia de Dirichlet} discreta, y reaparecer\'a
cuando analicemos el sobre-suavizado.

\subsection{La transformada de Fourier en grafos y los filtros espectrales}
\label{sec:graphs:gft}

Como $\Lap$ es sim\'etrico, admite una descomposici\'on en vectores propios
ortonormales $\Lap = U \Lambda U^\top$ con $\Lambda = \operatorname{diag}(\lambda_1,
\dots, \lambda_n)$ y $U = [u_1, \dots, u_n]$. Los vectores propios desempe\~nan el
papel de las exponenciales complejas del an\'alisis de Fourier cl\'asico
(v\'ease el \Cref{ch:grids}): el valor propio $\lambda_i$ es una noci\'on de
\emph{frecuencia}, con valores propios peque\~nos correspondientes a vectores
propios suaves y de variaci\'on lenta, y valores propios grandes a vectores que
oscilan r\'apidamente, como precisan el cociente de Rayleigh
$\lambda_i = u_i^\top \Lap\, u_i$ y la \Cref{prop:graphs:laplacian-props}.

\begin{definition}[Transformada de Fourier en grafos]
\label{def:graphs:gft}
La \emph{transformada de Fourier en grafos} de una se\~nal $f \in \R^n$ es
$\hat f = U^\top f$, con inversa $f = U \hat f = \sum_{i} \hat f_i\, u_i$. El
coeficiente $\hat f_i = \inner{f}{u_i}$ es la amplitud de $f$ en la frecuencia
$\lambda_i$.
\end{definition}

Por analog\'ia con el teorema de convoluci\'on cl\'asico ---la convoluci\'on en el
espacio es multiplicaci\'on en frecuencia--- se define la convoluci\'on en un
grafo como la multiplicaci\'on de los coeficientes de Fourier por un filtro.

\begin{definition}[Filtro espectral en grafos]
\label{def:graphs:spectral-filter}
Un \emph{filtro espectral} con funci\'on de transferencia
$g_\theta \colon \R \to \R$ act\'ua sobre una se\~nal $f$ mediante
\[
  g_\theta \star f
    = U\, g_\theta(\Lambda)\, U^\top f
    = U\, \operatorname{diag}\!\big(g_\theta(\lambda_1), \dots, g_\theta(\lambda_n)\big)\, U^\top f .
\]
\end{definition}

Esta es la base de las redes espectrales de grafos~\citep%
{bruna2014spectralnetworkslocallyconnected}, que aprenden directamente la funci\'on
de transferencia $g_\theta(\lambda)$. Por atractivo que sea, el enfoque espectral
ingenuo tiene tres inconvenientes. Primero, requiere la descomposici\'on
espectral completa de $\Lap$, un c\'omputo $O(n^3)$. Segundo, una
$g_\theta(\lambda)$ parametrizada libremente usa $O(n)$ par\'ametros y no est\'a
localizada en el espacio. Tercero ---y de forma m\'as fundamental--- la base de
vectores propios $U$ est\'a ligada a un grafo, de modo que un filtro aprendido en
un grafo no se transfiere a otro, lo que rompe el escenario inductivo en el que
los grafos var\'ian de un ejemplo a otro.

\subsection{Filtros polinomiales: el puente hacia los m\'etodos espaciales}
\label{sec:graphs:cheb}

Los tres inconvenientes se eliminan de golpe restringiendo la funci\'on de
transferencia a un \emph{polinomio} del laplaciano. Si $g_\theta(\lambda) =
\sum_{k=0}^{K} \theta_k\, \lambda^k$, entonces $g_\theta \star f = \big(\sum_k
\theta_k \Lap^k\big) f$, y los vectores propios se cancelan: no hace falta
descomposici\'on espectral alguna. Adem\'as $\Lap^k$ conecta solo nodos a
distancia de grafo a lo sumo $k$, de modo que un filtro polinomial de grado $K$
est\'a exactamente $K$-localizado ---es, por construcci\'on, una operaci\'on
espacial---. Este es el sentido preciso en que las visiones espectral y espacial
coinciden.

\citet{defferrard2017convolutionalneuralnetworksgraphs} lo hacen pr\'actico con
\emph{ChebNet}, que expande el filtro en polinomios de Chebyshev $T_k$ sobre el
espectro reescalado $\tilde\Lap = 2\Lap/\lambda_{\max} - I$:
\[
  g_\theta \star f
    = \sum_{k=0}^{K} \theta_k\, T_k(\tilde\Lap)\, f .
\]
La recurrencia de Chebyshev $T_k(x) = 2x\,T_{k-1}(x) - T_{k-2}(x)$ permite
calcular $T_k(\tilde\Lap) f$ mediante productos matriz--vector dispersos
repetidos, con coste total $O(K \abs{E})$ y solo $K+1$ par\'ametros, ninguno de
los cuales depende de la base de vectores propios. La red convolucional de grafos
de \citet{kipf2017semi} es el caso de primer orden $K = 1$ con un \'unico
par\'ametro compartido y una renormalizaci\'on de la matriz de adyacencia, lo que
da la capa $H^{(\ell+1)} = \sigma\!\big(\hat D^{-1/2} \hat A\, \hat D^{-1/2}
H^{(\ell)} W^{(\ell)}\big)$ con $\hat A = A + I$. Esta es claramente una capa de
paso de mensajes (\Cref{def:graphs:mpnn}) con agregaci\'on normalizada por grado,
lo que cierra el c\'irculo: la construcci\'on espectral, hecha local y eficiente
en par\'ametros, \emph{es} paso de mensajes espacial.

\section{Expresividad y la jerarqu\'ia de Weisfeiler--Leman}
\label{sec:graphs:wl}

Disponer de un mecanismo unificador invita a una pregunta aguda: \textquestiondown
cu\'an potente es? \textquestiondown Puede una red de paso de mensajes distinguir
dos grafos no isomorfos cualesquiera? La respuesta es no, y el l\'imite preciso lo
da un algoritmo combinatorio cl\'asico de la comprobaci\'on de isomorfismo de
grafos.

\subsection{El test de Weisfeiler--Leman}
\label{sec:graphs:wl-test}

\begin{definition}[Test de Weisfeiler--Leman de orden 1]
\label{def:graphs:wl}
Dado un grafo $G = (V, E)$, el algoritmo \emph{Weisfeiler--Leman de orden $1$}
($1$-\WL) asigna colores a los nodos y los refina iterativamente. Se inicializa
$c^{(0)}(v)$ id\'enticamente para todos los nodos (o a partir de etiquetas de
nodo, si existen). En cada paso,
\[
  c^{(t+1)}(v) = \HASH\!\Big(c^{(t)}(v),\,
      \{\!\{\, c^{(t)}(u) : u \in \Nbr(v) \,\}\!\}\Big),
\]
donde $\{\!\{\cdot\}\!\}$ es un multiconjunto y $\HASH$ es una aplicaci\'on
inyectiva que env\'ia cada par (color, multiconjunto de vecinos) distinto a un
color nuevo. El refinamiento se estabiliza en a lo sumo $n$ pasos; dos grafos son
\emph{$\WL$-equivalentes} cuando los multiconjuntos finales de colores (los
histogramas de colores) coinciden.
\end{definition}

El parecido estructural entre el paso de refinamiento y una capa de paso de
mensajes es inconfundible: un nodo actualiza su color a partir de su propio color
y del multiconjunto de los colores de sus vecinos, exactamente como un nodo
actualiza sus caracter\'isticas a partir de las propias y del agregado de los
mensajes de sus vecinos. La \Cref{fig:graphs:wl} ilustra el refinamiento.

\begin{figure}[h]
\centering
\begin{tikzpicture}[>=Stealth, font=\small,
    wlnode/.style={circle, draw, thick, minimum size=6mm, inner sep=0pt}]
\foreach \dx in {0,5.0,10.0}{
  \begin{scope}[shift={(\dx,0)}]
    \coordinate (n1) at (0,1.6);\coordinate (n2) at (1.8,1.6);
    \coordinate (n3) at (1.8,0);\coordinate (n4) at (0,0);
    \coordinate (n5) at (0.9,-1.1);\coordinate (n6) at (0.9,0.8);
    \draw[thick,gdlgray!50] (n1)--(n2) (n2)--(n3) (n3)--(n4) (n4)--(n1)
       (n3)--(n5) (n4)--(n5) (n1)--(n6) (n2)--(n6) (n3)--(n6) (n4)--(n6);
  \end{scope}
}
\begin{scope}
  \node at (0.9,2.6) {\bfseries (a) $c^{(0)}$};
  \foreach \p/\x/\y in {1/0/1.6,2/1.8/1.6,3/1.8/0,4/0/0,5/0.9/-1.1,6/0.9/0.8}
    \node[wlnode, fill=gdlgray!40] at (\x,\y) {};
\end{scope}
\draw[->, very thick, gdlgray!50] (2.7,0.6) -- node[above,font=\scriptsize]{HASH} (4.3,0.6);
\begin{scope}[shift={(5.0,0)}]
  \node at (0.9,2.6) {\bfseries (b) $c^{(1)}$};
  \foreach \p/\x/\y/\c in {1/0/1.6/{gdlblue!50},2/1.8/1.6/{gdlblue!50},3/1.8/0/{gdlred!50},4/0/0/{gdlred!50},5/0.9/-1.1/{gdlgreen!50},6/0.9/0.8/{gdlred!50}}
    \node[wlnode, fill=\c] at (\x,\y) {};
\end{scope}
\draw[->, very thick, gdlgray!50] (7.7,0.6) -- node[above,font=\scriptsize]{HASH} (9.3,0.6);
\begin{scope}[shift={(10.0,0)}]
  \node at (0.9,2.6) {\bfseries (c) $c^{(2)}$};
  \foreach \p/\x/\y/\c in {1/0/1.6/{gdlblue!50},2/1.8/1.6/{gdlblue!50},3/1.8/0/{gdlyellow!60},4/0/0/{gdlyellow!60},5/0.9/-1.1/{gdlgreen!50},6/0.9/0.8/{gdlred!50}}
    \node[wlnode, fill=\c] at (\x,\y) {};
\end{scope}
\end{tikzpicture}
\caption[Refinamiento de colores de Weisfeiler--Leman]{El refinamiento $1$-\WL{}
sobre un grafo de seis nodos. (a)~Coloraci\'on inicial uniforme. (b)~Tras un paso
los nodos se distinguen por su grado. (c)~Un paso adicional separa los dos nodos
de grado $4$ ($v_3, v_4$) de $v_6$, porque sus multiconjuntos de colores vecinos
difieren.}
\label{fig:graphs:wl}
\end{figure}

Este parecido se convierte en teorema gracias a \citet{xu2019how} y
\citet{morris2019weisfeiler}.

\begin{theorem}[El poder expresivo de las MPNN iguala a 1-\WL]
\label{thm:graphs:wl-mpnn}
Las redes neuronales de paso de mensajes est\'an acotadas superiormente por
$1$-\WL: si dos grafos son $\WL$-equivalentes, toda MPNN les asigna
representaciones id\'enticas. Adem\'as la cota es ajustada: existen MPNN (la Graph
Isomorphism Network de \citet{xu2019how}) que distinguen exactamente los grafos
que $1$-\WL{} distingue, lo que se logra usando funciones de agregaci\'on y
actualizaci\'on inyectivas.
\end{theorem}

\begin{proof}[Esbozo de demostraci\'on]
Para la cota superior, obs\'ervese que una capa MPNN calcula los estados de los
nodos a partir del estado propio y del multiconjunto de estados vecinos, la misma
informaci\'on que usa $1$-\WL; por inducci\'on sobre las capas, los grafos
$\WL$-equivalentes se aplican a representaciones iguales. Para la ajustez, se
construye un mensaje y una actualizaci\'on que realicen un \emph{hash} inyectivo.
La agregaci\'on por suma sobre un espacio de caracter\'isticas numerable, seguida
de un perceptr\'on multicapa, es inyectiva sobre multiconjuntos por el teorema de
aproximaci\'on universal para funciones de conjuntos~\citep{zaheer2017deepsets},
de modo que la MPNN puede reproducir exactamente cada paso de refinamiento.
\end{proof}

\begin{example}[Dos grafos que $1$-\WL{} no puede separar]
\label{ex:graphs:wl-fail}
Sea $G_1$ el ciclo de seis nodos $C_6$ y $G_2$ la uni\'on disjunta de dos
tri\'angulos $K_3 \sqcup K_3$. Ambos tienen seis nodos, seis aristas y todos sus
nodos de grado $2$. Bajo $1$-\WL{} todos los nodos empiezan con un color, ven el
multiconjunto $\{\!\{1,1\}\!\}$ y se estabilizan en un solo color; los histogramas
son id\'enticos, de modo que $G_1 \equiv_{\WL} G_2$. Sin embargo, los grafos no
son isomorfos: $C_6$ es conexo mientras que $K_3 \sqcup K_3$ tiene dos
componentes. Una MPNN est\'andar no puede, por tanto, distinguir un anillo
hexagonal de un par de tri\'angulos, un fallo instructivo cuando la conectividad
global importa ---por ejemplo, al distinguir ciertos is\'omeros estructurales en
qu\'imica---.
\end{example}

\subsection{Sobre-suavizado}
\label{sec:graphs:oversmoothing}

El \Cref{thm:graphs:wl-mpnn} acota lo que una red puede distinguir, pero no dice
nada sobre lo que sucede al apilar muchas capas. Emp\'iricamente, las MPNN
profundas se degradan: las representaciones de los nodos convergen a un valor
com\'un y pierden todo poder discriminativo. Este es el \emph{sobre-suavizado}
(\emph{oversmoothing})~\citep{li2018deeper}, y admite una explicaci\'on espectral
n\'itida a trav\'es de la energ\'ia de Dirichlet.

\begin{definition}[Energ\'ia de Dirichlet]
\label{def:graphs:dirichlet}
Para caracter\'isticas de nodos $H \in \R^{n \times d}$, la \emph{energ\'ia de
Dirichlet} es
\[
  E(H) = \tr(H^\top L\, H)
       = \sum_{(i,j) \in E} W_{ij}\,\norm{h_i - h_j}^2 ,
\]
que se anula si y solo si $H$ es constante en cada componente conexa.
\end{definition}

Un paso de paso de mensajes lineal con agregaci\'on por media es un paso de
difusi\'on en el grafo: si $h_v^{(\ell+1)} = (1-\alpha) h_v^{(\ell)} + \alpha
\sum_{u \in \Nbr(v)} h_u^{(\ell)}/\deg(v)$, entonces en forma matricial
$H^{(\ell+1)} = (I - \alpha \tilde L)\, H^{(\ell)}$ con el laplaciano de paseo
aleatorio $\tilde L = D^{-1} L$. Esta es la discretizaci\'on de Euler expl\'icita
de la ecuaci\'on del calor $\partial_t H = -\alpha \tilde L H$ en el grafo, e
iterarla suaviza la se\~nal.

\begin{theorem}[Sobre-suavizado exponencial]
\label{thm:graphs:oversmoothing}
Sea $G$ conexo con valores propios del laplaciano normalizado
$0 = \lambda_1 < \lambda_2 \le \cdots \le \lambda_n$ y sea
$0 < \alpha < 1/\lambda_n$. Entonces tras $\ell$ capas de difusi\'on lineal
\[
  E\big(H^{(\ell)}\big) \le (1 - \alpha \lambda_2)^{2\ell}\, E\big(H^{(0)}\big),
\]
de modo que $E(H^{(\ell)}) \to 0$ exponencialmente: todas las representaciones de
los nodos colapsan a un vector com\'un.
\end{theorem}

\begin{proof}[Esbozo de demostraci\'on]
P\'asese al laplaciano normalizado sim\'etrico $L_{\mathrm{sym}}$, que comparte
espectro con $\tilde L$, y exp\'andanse las caracter\'isticas en su base de
vectores propios. El operador de difusi\'on escala el coeficiente de $u_k$ por
$(1 - \alpha\lambda_k)$, de modo que tras $\ell$ pasos la energ\'ia es
$\sum_{k \ge 2} \lambda_k (1-\alpha\lambda_k)^{2\ell} \abs{\hat f_k}^2$ (el
t\'ermino $k=1$ se anula porque $\lambda_1 = 0$). Como $\alpha < 1/\lambda_n$,
cada factor cumple $\abs{1-\alpha\lambda_k} \le 1 - \alpha\lambda_2$ para
$k \ge 2$, lo que da la cota.
\end{proof}

La tasa la gobierna el \emph{salto espectral} $\lambda_2$: los grafos bien
conectados (con $\lambda_2$ grande, como los expansores) se homogeneizan m\'as
deprisa. Esta es la sombra discreta de la difusi\'on del calor en una variedad
riemanniana, donde el salto espectral del operador de Laplace--Beltrami controla
la tasa de mezclado. Las conexiones residuales, la normalizaci\'on y la difusi\'on
en haces~\citep{bodnar2022neural} figuran entre los remedios.

\subsection{Sobre-compresi\'on y tests de orden superior}
\label{sec:graphs:oversquashing}

Si el sobre-suavizado es la p\'erdida de se\~nal de alta frecuencia al crecer la
profundidad, el fallo complementario es la \emph{sobre-compresi\'on}
(\emph{oversquashing}): la incapacidad de propagar informaci\'on a larga
distancia, porque los mensajes de un vecindario a $\ell$ saltos que crece de forma
exponencial deben comprimirse en un vector de tama\~no fijo. La sensibilidad de un
nodo a una entrada lejana, $\norm{\partial h_v^{(\ell)} / \partial h_u^{(0)}}$,
est\'a acotada por $(c_M c_U)^\ell\, [\hat A^\ell]_{vu}$ con $\hat A = D^{-1} A$ y
constantes de Lipschitz $c_M, c_U$ de las funciones de mensaje y actualizaci\'on;
a trav\'es de los cuellos de botella topol\'ogicos la entrada $[\hat A^\ell]_{vu}$
decae de forma abrupta. \citet{topping2022understanding} ligan estos cuellos de
botella a la curvatura de Ricci discreta negativa en las aristas y proponen el
\emph{recableado} (\emph{rewiring}) del grafo para aliviarlos ---un ejemplo
limpio de geometr\'ia diferencial que prescribe una cura para una patolog\'ia del
aprendizaje, en paralelo a c\'omo la curvatura seccional negativa hace divergir
las geod\'esicas (\Cref{ch:geodesics})---.

Por \'ultimo, la barrera $1$-\WL{} puede elevarse ascendiendo por una jerarqu\'ia
de tests m\'as potentes.

\begin{definition}[Test de Weisfeiler--Leman de orden $k$]
\label{def:graphs:k-wl}
El test \emph{$k$-\WL} colorea $k$-tuplas de nodos en lugar de nodos individuales,
refinando el color de cada tupla a partir de los colores de todas las tuplas
obtenidas al reemplazar una de sus entradas por un nodo
arbitrario~\citep{morris2019weisfeiler}.
\end{definition}

\begin{proposition}[Jerarqu\'ia estricta $k$-\WL]
\label{prop:graphs:k-wl}
Los tests forman una jerarqu\'ia estrictamente creciente en poder de
distinci\'on,
\[
  1\text{-}\WL \subsetneq 2\text{-}\WL \subsetneq 3\text{-}\WL \subsetneq \cdots,
\]
distinguiendo cada nivel grafos que el anterior no puede, a un coste de $O(n^k)$
de memoria por iteraci\'on.
\end{proposition}

El coste exponencial de $k$-\WL{} motiva v\'ias m\'as baratas para superar el techo
de $1$-\WL: codificaciones estructurales y posicionales (por ejemplo, vectores
propios del laplaciano de la \Cref{def:graphs:gft}) que rompen las simetr\'ias de
los nodos, \emph{transformers} de grafos espectrales~\citep{kreuzer2021spectral} y
el paso de mensajes en dominios topol\'ogicos de orden superior (complejos
simpliciales y celulares), que se abordan en los
\Cref{ch:gauges,ch:applications}.

\medskip
Las tres limitaciones ---sobre-suavizado, sobre-compresi\'on y barrera
$\WL$--- son facetas de un mismo hecho: el paso de mensajes es fundamentalmente
\emph{local}, y mueve informaci\'on un salto por capa. Demasiadas capas destruyen
la se\~nal; demasiado pocas impiden la comunicaci\'on a larga distancia; y el
c\'omputo puramente local no puede ver ciertos patrones globales. Las herramientas
geom\'etricas desarrolladas aqu\'i ---an\'alisis espectral, curvatura discreta,
topolog\'ia de orden superior--- no solo diagnostican estos l\'imites, sino que
prescriben remedios concretos, ejemplificando la tesis de este libro de que la
geometr\'ia es el principio organizador de las arquitecturas eficaces.
